% ============================================================
%  Глава 10 — Физический симулятор: кинематика, сенсоры, камера
%  Файл: compute_node/simulator.py
% ============================================================
\chapter{Физический симулятор: кинематика, сенсоры, камера}
\label{ch:simulator}

\begin{filebox}[title={Файл исходного кода}]
\filepath{compute\_node/simulator.py} \quad (строки 282--588, Python)
\end{filebox}

% ── Кинематика ─────────────────────────────────────────────────
\section{Кинематика робота (SimRobot)}

\begin{filebox}[title={\filepath{simulator.py}, строки 373--397 --- tick()}]
\begin{lstlisting}[style=pythonstyle, numbers=left, firstnumber=373]
def tick(self, dt: float):
    self._prev_v_linear = self.v_linear
    prev_x, prev_y = self.x, self.y
    # Integrate velocities (unicycle model)
    self.x += self.v_linear * math.cos(self.theta) * dt
    self.y += self.v_linear * math.sin(self.theta) * dt
    self.theta += self.v_angular * dt
    # Normalise theta
    self.theta = math.atan2(math.sin(self.theta), math.cos(self.theta))
    # Wall collision -- clamp inside arena
    margin = ROBOT_RADIUS
    self.x = max(margin, min(self.arena.width - margin, self.x))
    self.y = max(margin, min(self.arena.height - margin, self.y))
    # Forbidden zone collision -- push back
    for z in self.arena.forbidden_zones:
        zx1 = z['x1'] - margin
        zy1 = z['y1'] - margin
        zx2 = z['x2'] + margin
        zy2 = z['y2'] + margin
        if zx1 <= self.x <= zx2 and zy1 <= self.y <= zy2:
            self.x, self.y = prev_x, prev_y
            self.v_linear = 0.0
            break
\end{lstlisting}
\end{filebox}

\subsection{Дифференциальные уравнения кинематики}

Робот моделируется как \textbf{unicycle} --- одноколейная модель
с линейной и угловой скоростями:

\begin{align}
  \dot{x}      &= v \cos\theta \label{eq:sim_dx}\\
  \dot{y}      &= v \sin\theta \label{eq:sim_dy}\\
  \dot{\theta}  &= \omega       \label{eq:sim_dtheta}
\end{align}

Дискретизация методом Эйлера ($\Delta t = 0.05\,\text{с}$):

\begin{equation}
  \boxed{
  \begin{pmatrix} x_{k+1} \\ y_{k+1} \\ \theta_{k+1} \end{pmatrix}
  = \begin{pmatrix} x_k \\ y_k \\ \theta_k \end{pmatrix}
  + \Delta t \begin{pmatrix} v\cos\theta_k \\ v\sin\theta_k \\ \omega \end{pmatrix}
  }
  \label{eq:sim_euler}
\end{equation}

\subsection{Нормализация угла}

Для предотвращения накопления $\theta > 2\pi$:

\begin{equation}
  \theta \leftarrow \mathrm{atan2}\!\bigl(\sin\theta,\, \cos\theta\bigr)
  \label{eq:theta_norm}
\end{equation}

\subsection{Коллизия со стенами}

Ограничение позиции внутри арены:

\begin{equation}
  x \leftarrow \mathrm{clamp}(x,\; r_{\text{robot}},\; W - r_{\text{robot}}), \quad
  y \leftarrow \mathrm{clamp}(y,\; r_{\text{robot}},\; H - r_{\text{robot}})
  \label{eq:wall_collision}
\end{equation}

% ── Сенсоры ────────────────────────────────────────────────────
\section{Ультразвуковой сенсор (SimSensors)}

\begin{filebox}[title={\filepath{simulator.py}, строки 417--464 --- \_update\_ultrasonic()}]
\begin{lstlisting}[style=pythonstyle, numbers=left, firstnumber=417]
def _update_ultrasonic(self, robot, arena):
    """Ray-cast forward from robot."""
    rx, ry = robot.x, robot.y
    dx = math.cos(robot.theta)
    dy = math.sin(robot.theta)
    best = ULTRASONIC_MAX

    # Check walls (ray-plane intersection)
    if dx > 0:
        t = (arena.width - rx) / dx     # right wall
        if ULTRASONIC_MIN < t < best: best = t
    if dx < 0:
        t = -rx / dx                    # left wall
        if ULTRASONIC_MIN < t < best: best = t
    if dy > 0:
        t = (arena.height - ry) / dy    # top wall
        if ULTRASONIC_MIN < t < best: best = t
    if dy < 0:
        t = -ry / dy                    # bottom wall
        if ULTRASONIC_MIN < t < best: best = t

    # Check balls (ray-sphere intersection)
    for ball in arena.balls:
        if ball['grabbed']: continue
        to_ball_x = ball['x'] - rx
        to_ball_y = ball['y'] - ry
        proj = to_ball_x * dx + to_ball_y * dy   # projection
        perp = abs(to_ball_x * dy - to_ball_y * dx)  # perpendicular
        if proj > ULTRASONIC_MIN and perp < ball['radius'] + 0.05:
            if proj < best: best = proj

    # Gaussian noise
    noise = random.gauss(0, 0.005)
    self.range_m = max(ULTRASONIC_MIN, min(ULTRASONIC_MAX, best + noise))
\end{lstlisting}
\end{filebox}

\subsection{Пересечение луч--плоскость (стены)}

Луч из точки $(r_x, r_y)$ по направлению $(\cos\theta, \sin\theta)$:

\begin{equation}
  \vect{p}(t) = \begin{pmatrix} r_x \\ r_y \end{pmatrix}
  + t \begin{pmatrix} \cos\theta \\ \sin\theta \end{pmatrix}, \quad t > 0
  \label{eq:ray_def}
\end{equation}

Расстояние до стены $x = W$ (при $\cos\theta > 0$):

\begin{equation}
  t = \frac{W - r_x}{\cos\theta}
  \label{eq:ray_wall}
\end{equation}

\subsection{Пересечение луч--сфера (шарики)}

Для каждого шарика с центром $(b_x, b_y)$:

\textbf{Проекция} вектора ``к шарику'' на направление луча:

\begin{equation}
  t_{\text{proj}} = (\vect{b} - \vect{r}) \cdot \hat{\vect{d}}
  = (b_x - r_x)\cos\theta + (b_y - r_y)\sin\theta
  \label{eq:ball_proj}
\end{equation}

\textbf{Перпендикулярное расстояние} от луча до центра шарика:

\begin{equation}
  d_{\perp} = |(b_x - r_x)\sin\theta - (b_y - r_y)\cos\theta|
  \label{eq:ball_perp}
\end{equation}

Условие обнаружения: $d_\perp < r_{\text{ball}} + 0.05$ (конус ультразвука).

\subsection{Шумовая модель}

\begin{equation}
  r_{\text{meas}} = \mathrm{clamp}(r_{\text{true}} + \mathcal{N}(0, 0.005^2),\;
  r_{\min},\; r_{\max})
  \label{eq:ultrasonic_noise}
\end{equation}

% ── IMU ────────────────────────────────────────────────────────
\section{Симуляция IMU}

\begin{filebox}[title={\filepath{simulator.py}, строки 466--472 --- \_update\_imu()}]
\begin{lstlisting}[style=pythonstyle, numbers=left, firstnumber=466]
def _update_imu(self, robot):
    noise = random.gauss(0, 0.3)
    self.imu_yaw = math.degrees(robot.theta) + noise
    self.imu_pitch = random.gauss(0, 0.2)
    self.imu_roll  = random.gauss(0, 0.2)
    self.imu_gyro_z = robot.v_angular
    self.accel_x = (robot.v_linear - robot._prev_v_linear) / SIM_DT
\end{lstlisting}
\end{filebox}

Линейное ускорение из конечной разности скорости:

\begin{equation}
  a_x = \frac{v_k - v_{k-1}}{\Delta t}
  \label{eq:sim_accel}
\end{equation}

% ── Камера ─────────────────────────────────────────────────────
\section{Модель перспективной камеры (SimDetector)}

\begin{filebox}[title={\filepath{simulator.py}, строки 510--543 --- \_render\_camera()}]
\begin{lstlisting}[style=pythonstyle, numbers=left, firstnumber=514]
# Vector from robot to ball
dx = bx - robot.x
dy = by - robot.y
dist = math.sqrt(dx*dx + dy*dy)

# Angle to ball relative to robot heading
angle = math.atan2(dy, dx) - robot.theta
angle = math.atan2(math.sin(angle), math.cos(angle))

if abs(angle) > CAM_FOV / 2:
    continue   # outside field of view

# Project to screen
screen_x = int(CAM_W/2 + (angle / (CAM_FOV/2)) * (CAM_W/2))

# Apparent size (pinhole camera model)
apparent_px = int(FOCAL_LENGTH_PX * BALL_DIAMETER_M / dist)

# Vertical position (depth cue)
screen_y = horizon_y + int((1.0 - 0.03/max(0.1, dist)) *
                            (CAM_H - horizon_y) * 0.7)

# Confidence proportional to distance
conf = max(0.5, min(0.99, 1.0 - dist / 3.0))
\end{lstlisting}
\end{filebox}

\subsection{Проверка видимости (Field of View)}

Угол на объект относительно курса робота:

\begin{equation}
  \alpha = \mathrm{atan2}(b_y - r_y,\, b_x - r_x) - \theta_{\text{robot}}
  \label{eq:fov_angle}
\end{equation}

Объект виден если: $|\alpha| \leq \text{FOV}/2 = 30^\circ$.

\subsection{Горизонтальная проекция на экран}

\begin{equation}
  \boxed{s_x = \frac{W}{2} + \frac{\alpha}{\text{FOV}/2} \cdot \frac{W}{2}}
  \label{eq:screen_x}
\end{equation}

Это линейная проекция: угол $\alpha = 0$ --- центр кадра, $\pm \text{FOV}/2$ --- края.

\subsection{Модель pinhole-камеры (размер объекта)}

Видимый размер объекта на сенсоре:

\begin{equation}
  \boxed{s_{\text{px}} = \frac{f \cdot D_{\text{real}}}{d}}
  \label{eq:pinhole}
\end{equation}

где $f = 500\,\text{пкс}$ --- фокусное расстояние, $D_{\text{real}} = 0.04\,\text{м}$ ---
диаметр шарика, $d$ --- расстояние до объекта.

\begin{notebox}
Это обратная зависимость: при удвоении расстояния видимый размер уменьшается вдвое.
Тот же принцип используется в \texttt{follow\_me\_node.py}
(глава~\ref{ch:follow_me}) для оценки дистанции до человека.
\end{notebox}

\subsection{Вертикальная позиция (глубинная подсказка)}

Объекты на полу: чем ближе к роботу, тем ниже на экране:

\begin{equation}
  s_y = h_{\text{горизонт}} + \left(1 - \frac{0.03}{\max(0.1, d)}\right) \cdot 0.7(H - h_{\text{горизонт}})
  \label{eq:screen_y}
\end{equation}

\subsection{Модель доверия (confidence)}

\begin{equation}
  \text{conf} = \mathrm{clamp}\!\left(1 - \frac{d}{3.0},\; 0.5,\; 0.99\right)
  \label{eq:sim_conf}
\end{equation}

% ── Таблица ────────────────────────────────────────────────────
\section{Параметры симулятора}

\begin{center}
\begin{tabular}{lll}
\toprule
\textbf{Параметр} & \textbf{Значение} & \textbf{Описание}\\
\midrule
\texttt{ARENA\_W, H}        & $3.0 \times 3.0\,\text{м}$ & Размер арены\\
\texttt{SIM\_DT}            & $0.05\,\text{с}$ & Шаг интегрирования (20~Гц)\\
\texttt{MAX\_LINEAR}        & $0.3\,\text{м/с}$ & Макс. линейная скорость\\
\texttt{MAX\_ANGULAR}       & $2.0\,\text{рад/с}$ & Макс. угловая скорость\\
\texttt{ROBOT\_RADIUS}      & $0.12\,\text{м}$ & Радиус робота\\
\texttt{BALL\_DIAMETER\_M}  & $0.04\,\text{м}$ & Диаметр шарика\\
\texttt{FOCAL\_LENGTH\_PX}  & $500\,\text{пкс}$ & Фокусное расстояние камеры\\
\texttt{CAM\_FOV}           & $60^\circ$ & Поле зрения камеры\\
\texttt{CAM\_W, H}          & $640 \times 480\,\text{пкс}$ & Разрешение камеры\\
\texttt{ULTRASONIC\_MAX}    & $2.0\,\text{м}$ & Макс. дальность УЗ\\
\bottomrule
\end{tabular}
\end{center}
